\documentclass[a4paper,12pt]{article}
\usepackage[utf8]{inputenc}
\usepackage[T1]{fontenc}
\usepackage[ngerman]{babel}
\usepackage{amsmath}
\usepackage{amssymb}
\usepackage{enumitem}
\usepackage{geometry}
\geometry{a4paper, margin=2.5cm}
\usepackage{parskip}

\title{Einsendeaufgaben -- Aufgabe 4.3\\[0.5em]
\large Modul 61112: Lineare Algebra}
\author{}
\date{}

\begin{document}
\maketitle

\section*{Aufgabe 4.3}

Anwendung des Adjunktensatzes:
\begin{align*}
  \mathrm{Adj}(B)\, \mathrm{Adj}(A)\, AB
    &= \mathrm{Adj}(B)\, \det(A)\, I_n\, B && (\text{Adjunktensatz}) \\
    &= \mathrm{Adj}(B)\, B\, \det(A)\, I_n && (\text{MG 2.4.1c}) \\
    &= \det(B)\, I_n\, \det(A)\, I_n && (\text{Adjunktensatz}) \\
    &= \det(B)\, \det(A)\, I_n && (\text{MG 2.4.1c}) \\
    &= \det(A)\, \det(B)\, I_n && (\text{Eigenschaft Körper } K) \\
    &= \det(AB)\, I_n && (\text{Determinantenmultiplikationssatz}) \\
    &= \mathrm{Adj}(AB)\, AB && (\text{Adjunktensatz})
\end{align*}

Da $A$ und $B$ invertierbar sind und $\mathrm{Adj}(B)\, \mathrm{Adj}(A)\, AB = \mathrm{Adj}(AB)\, AB$ gilt, folgt
\[
  \mathrm{Adj}(B)\, \mathrm{Adj}(A) = \mathrm{Adj}(AB).
\]

\end{document}
